\begin{abstract}
    {\bf Background:} 
    Agents for computer use (ACUs) are systems that execute complex tasks on digital devices -- such as personal computers or mobile phones -- given instructions in natural language. These agents automate tasks by controlling software through low-level actions like mouse clicks and touchscreen gestures. However, despite rapid progress, ACUs are not yet mature for everyday use.

    
    {\bf Objectives:}
    This survey examines the current state-of-the-art, identifies trends, and points out research gaps in the development of practical ACUs. The goal is to provide a comprehensive review and analysis that helps advance general-purpose, robust, and scalable agents for real-world computer use.

    
    {\bf Methods:}
    We introduce a multifaceted taxonomy of ACUs across three dimensions: (\uproman{1}) the \emph{domain perspective}, characterizing the contexts in which agents operate; (\uproman{2}) the \emph{interaction perspective}, describing observation modalities (e.g., screenshots, HTML) and action modalities (e.g., mouse, keyboard, code execution); and (\uproman{3}) the \emph{agent perspective}, detailing how agents perceive, reason, and learn. We review $\numagents$ original research papers about ACUs and $\numdatasets$ relevant datasets, covering both foundation model-based and specialized approaches.  
    
    {\bf Results:}
    Our taxonomy comprehensively structures state-of-the-art approaches and establishes the groundwork for guiding future ACU research. We found that the field is transitioning from specialized agents toward foundation-model-based agents, a shift from text to image-based observation space, and an increasing adoption of behavior cloning methodologies. Furthermore, we identify six key research gaps: insufficient generalization, inefficient learning, limited planning, low task complexity in benchmarks, non-standardized evaluation, and a disconnect between research and practical conditions.
    
    {\bf Conclusions:}
    To continue rapid improvements in the field, we recommend focusing on: (a) vision-based observations and low-level control to enhance generalization; (b) adaptive learning beyond static prompting; (c) effective planning and reasoning capabilities; (d) realistic, high-complexity benchmarks; (e) standardized evaluation criteria based on task success; and (f) aligning agent design with real-world deployment constraints. Collectively, our findings and proposed directions help develop more general-purpose agents for everyday digital tasks.
\end{abstract}